\documentclass[12pt, a4paper]{report}

% ── Encodage & langue ───────────────────────────────────────
\usepackage[utf8]{inputenc}
\usepackage[T1]{fontenc}
\usepackage[french]{babel}

% ── Mise en page ────────────────────────────────────────────
\usepackage[
  top=2.5cm, bottom=2.5cm,
  left=3cm,  right=2.5cm,
  headheight=15pt
]{geometry}
\usepackage{setspace}
\onehalfspacing                     % interligne 1,5

% ── Polices & typographie ───────────────────────────────────
\usepackage{lmodern}
\usepackage{microtype}

% ── En-têtes / pieds de page ────────────────────────────────
\usepackage{fancyhdr}
\pagestyle{fancy}
\fancyhf{}
\fancyhead[L]{\small\leftmark}
\fancyhead[R]{\small\thepage}
\fancyfoot[C]{\small École Supérieure de Technologie — Fkih Ben Salah}
\renewcommand{\headrulewidth}{0.4pt}
\renewcommand{\footrulewidth}{0.4pt}

% ── Couleurs ────────────────────────────────────────────────
\usepackage{xcolor}
\definecolor{estblue}{RGB}{0, 70, 127}      % bleu institutionnel
\definecolor{lightgray}{RGB}{245, 245, 245}

% ── Graphiques & figures ────────────────────────────────────
\usepackage{graphicx}
\usepackage{float}
\usepackage{caption}
\usepackage{subcaption}
\graphicspath{{images/}}            % dossier des images

% ── Tableaux ────────────────────────────────────────────────
\usepackage{booktabs}
\usepackage{array}
\usepackage{tabularx}
\usepackage{longtable}
\usepackage{multirow}

% ── Listes ──────────────────────────────────────────────────
\usepackage{enumitem}

% ── Mathématiques ───────────────────────────────────────────
\usepackage{amsmath, amssymb}

% ── Code source ─────────────────────────────────────────────
\usepackage{listings}
\usepackage{lstautogobble}
\lstset{
  basicstyle=\ttfamily\footnotesize,
  keywordstyle=\color{estblue}\bfseries,
  commentstyle=\color{gray}\itshape,
  stringstyle=\color{orange},
  numbers=left,
  numberstyle=\tiny\color{gray},
  numbersep=8pt,
  breaklines=true,
  breakatwhitespace=true,
  frame=single,
  backgroundcolor=\color{lightgray},
  captionpos=b,
  showstringspaces=false,
  tabsize=2,
  autogobble=true,
}
\lstdefinelanguage{SQL}{
  keywords={CREATE,TABLE,PRIMARY,KEY,FOREIGN,REFERENCES,
            NOT,NULL,DEFAULT,UNIQUE,AUTO_INCREMENT,
            INSERT,SELECT,UPDATE,DELETE,USE,DATABASE,IF,EXISTS},
  sensitive=false,
  morecomment=[l]{--},
  morestring=[b]',
}

% ── Hyperliens ──────────────────────────────────────────────
\usepackage[
  colorlinks=true,
  linkcolor=estblue,
  urlcolor=estblue,
  citecolor=estblue,
  bookmarks=true,
  bookmarksopen=true
]{hyperref}

% ── Références ──────────────────────────────────────────────
\usepackage[backend=biber, style=numeric, sorting=none]{biblatex}
\addbibresource{references.bib}     % fichier BibTeX

% ── Numérotation des sections dans les chapitres ────────────
\setcounter{secnumdepth}{3}
\setcounter{tocdepth}{3}

% ── Style des titres de chapitres ───────────────────────────
\usepackage{titlesec}
\titleformat{\chapter}[display]
  {\normalfont\huge\bfseries\color{estblue}}
  {\chaptertitlename\ \thechapter}{20pt}{\Huge}
\titlespacing*{\chapter}{0pt}{30pt}{20pt}

% ── Boîtes colorées (optionnel) ─────────────────────────────
\usepackage{tcolorbox}
\tcbuselibrary{skins}
\newtcolorbox{resumebox}{
  colback=lightgray, colframe=estblue,
  title=Résumé du chapitre, fonttitle=\bfseries,
  arc=4pt
}

% ── Abréviations / Acronymes ────────────────────────────────
\usepackage[acronym, toc]{glossaries}
\makeglossaries
% Exemple : \newacronym{api}{API}{Application Programming Interface}

% ── Méta-données PDF ────────────────────────────────────────
\hypersetup{
  pdftitle={Application web de réservation intelligente — Grand Palais},
  pdfauthor={Mohsine DAHNAOUI, Anas SABIQ, Ilyas QUIOAMI},
  pdfsubject={Génie Informatique — EST-FBS},
}

% ============================================================
\begin{document}
% ============================================================

% ── PAGE DE GARDE ───────────────────────────────────────────
\begin{titlepage}
  \centering

  % Logos
  \begin{minipage}{0.25\textwidth}
    
\includegraphics[width=\linewidth]{logo_est.png}   % logo EST-FBS
  \end{minipage}
  \hfill
  \begin{minipage}{0.55\textwidth}
    \centering
    {\small\textbf{Ministère de l'Enseignement Supérieur\\
    de la Recherche Scientifique et de l'Innovation\\[2pt]
    Université Sultan Moulay Slimane\\
    L'École Supérieure de Technologie\\
    Fkih Ben Salah}}
  \end{minipage}
  \hfill
  \begin{minipage}{0.12\textwidth}
    
\includegraphics[width=\linewidth]{logo_usms.png}  % logo USMS
  \end{minipage}

  \vspace{1.5cm}
  \rule{\linewidth}{1pt}\\[0.4cm]
  {\Large\textbf{Rapport de Projet de Fin d'Études}}\\[0.2cm]
  {\large Présenté en vue de l'obtention du}\\[0.1cm]
  {\large\textbf{Diplôme Universitaire de Technologie}}\\[0.1cm]
  {\large Filière : \textbf{Génie Informatique}}\\[0.2cm]
  \rule{\linewidth}{1pt}

  \vspace{1cm}

  % Auteurs
  {\large Présenté et Soutenu par :}\\[0.3cm]
  {\Large\textbf{Mr. Mohsine DAHNAOUI, Mr. Anas SABIQ et Mr. Ilyas QUIOAMI}}

  \vspace{0.8cm}

  % Encadrant
  {\large Encadré par :}\\[0.2cm]
  {\large\textbf{Pr. Khalid QBOUCHE}}

  \vspace{1cm}

  % Titre du projet
  \begin{tcolorbox}[
    colback=white, colframe=estblue,
    boxrule=1.5pt, arc=6pt,
    width=0.85\linewidth
  ]
    \centering\Large\bfseries
    Application web de réservation intelligente pour un hôtel\\[4pt]
    {\normalsize\normalfont\textit{Grand Palais — Système de Gestion Hôtelière Moderne}}
  \end{tcolorbox}

  \vfill

  % Jury
  {\large Soutenu le \textbf{JJ/MM/AAAA} devant la commission d'examen composée de :}\\[0.5cm]
  \begin{tabular}{ll}
    \textbf{Pr. Prénom NOM} & Encadrant \\
    \textbf{Pr. Prénom NOM} & Examinateur \\
    \textbf{Pr. Prénom NOM} & Président \\
  \end{tabular}

  \vspace{0.8cm}
  {\large Année universitaire \textbf{2024-2025}}

  \vspace{0.5cm}
  \rule{\linewidth}{0.4pt}\\
  {\footnotesize École Supérieure de Technologie — Fkih Ben Salah\\
  Hay Tighnari, Route Nationale N°11, 23200, B.P : 336 Fkih Ben Salah\\
  Tél : (+212) 05.23.43.46.66 \quad Fax : (+212) 05.23.43.49.99\\
  Email : estfbs@usms.ma \quad Site Web : \url{http://estfbs.usms.ac.ma/}}
\end{titlepage}

% ── PAGES LIMINAIRES (numérotation romaine) ──────────────────
\pagenumbering{roman}
\setcounter{page}{1}

% ── DÉDICACE ────────────────────────────────────────────────
\chapter*{Dédicace}
\addcontentsline{toc}{chapter}{Dédicace}
\thispagestyle{plain}

\begin{flushright}
  \itshape
  À nos parents,\\
  pour leur soutien inconditionnel et leurs sacrifices\\
  qui ont rendu ce parcours possible.\\[1em]
  À nos professeurs de l'École Supérieure de Technologie,\\
  qui nous ont transmis les connaissances et la passion\\
  pour l'informatique.\\[1em]
  À notre encadrant Pr. Khalid QBOUCHE,\\
  pour ses conseils précieux et son accompagnement\\
  tout au long de ce projet.\\[1em]
  À tous ceux qui croient en l'innovation au service de l'éducation.
\end{flushright}

\newpage

% ── REMERCIEMENTS ───────────────────────────────────────────
\chapter*{Remerciements}
\addcontentsline{toc}{chapter}{Remerciements}
\thispagestyle{plain}

Nous souhaitons adresser nos sincères remerciements à l'École Supérieure de Technologie
Fkih Ben Salah pour nous avoir donné l'opportunité de mener à bien ce projet de fin d'études.

Nos remerciements s'adressent tout particulièrement à \textbf{Pr. Khalid QBOUCHE},
notre encadrant pédagogique, pour son accompagnement rigoureux, ses conseils précieux et
sa disponibilité constante durant toute la durée de ce travail.

Nous exprimons également notre reconnaissance à l'ensemble de l'équipe pédagogique du
département Mathématiques et Informatique pour la qualité de leur enseignement et leurs
orientations expertes qui ont considérablement enrichi notre approche.

Nous remercions chaleureusement tous ceux qui ont contribué à faire de cette expérience
un parcours enrichissant.

\newpage

% ── RÉSUMÉ (français) ───────────────────────────────────────
\chapter*{Résumé}
\addcontentsline{toc}{chapter}{Résumé}
\thispagestyle{plain}

Le présent rapport détaille le développement de \textbf{Grand Palais}, une application web
de réservation hôtelière intelligente réalisée dans le cadre de notre projet de fin
d'études au département Mathématiques et Informatique de l'École Supérieure de Technologie
Fkih Ben Salah. Ce projet répond à la problématique de digitalisation et de modernisation
de la gestion opérationnelle des établissements hôteliers.

L'objectif principal consistait à concevoir et implémenter une plateforme complète couvrant
l'ensemble du cycle hôtelier : la réservation en ligne avec visualisation 3D interactive
des chambres, le paiement sécurisé via Stripe, la gestion des arrivées et départs
(check-in/check-out), et l'assistance aux clients par un concierge intelligent basé
sur l'IA. Le système implique trois acteurs distincts : le Client, le Réceptionniste
et l'Administrateur.

La démarche adoptée s'appuie sur une analyse fonctionnelle approfondie suivie d'une
conception UML rigoureuse, puis d'une implémentation utilisant Laravel~12 (PHP~8.2)
pour le backend et Vue.js~3 pour le frontend. L'architecture MVC retenue garantit la
séparation des responsabilités et facilite la maintenance évolutive.

Les résultats obtenus démontrent une solution robuste, sécurisée et innovante répondant
efficacement aux besoins de digitalisation du secteur hôtelier.

\bigskip
\noindent\textbf{Mots-clés :} Gestion hôtelière, Laravel, Vue.js, Réservation en ligne, Visualisation 3D, Intelligence artificielle, Paiement en ligne.

\newpage

% ── ABSTRACT (anglais) ──────────────────────────────────────
\chapter*{Abstract}
\addcontentsline{toc}{chapter}{Abstract}
\thispagestyle{plain}

This report details the development of \textbf{Grand Palais}, an intelligent hotel
reservation web application carried out as part of our final year project at the
Mathematics and Computer Science department of the Higher School of Technology Fkih
Ben Salah. This project addresses the digitalization and modernization of operational
management in the hospitality industry.

The main objective was to design and implement a complete platform covering the full
hotel lifecycle: online room reservation with interactive 3D visualization, secure
payment via Stripe, check-in and check-out management, and AI-powered concierge
assistance via WhatsApp. The system involves three distinct actors: the Client,
the Receptionist, and the Administrator.

The approach is based on thorough functional analysis followed by rigorous UML design,
then implementation using Laravel~12 (PHP~8.2) for the backend and Vue.js~3 with Vite
for the frontend, connected to a MySQL database.

\bigskip
\noindent\textbf{Keywords:} Hotel management, Laravel, Vue.js, Online booking, 3D visualization, Artificial intelligence, Online payment.

\newpage

% ── LISTE DES ACRONYMES ─────────────────────────────────────
\chapter*{Liste des Acronymes}
\addcontentsline{toc}{chapter}{Liste des Acronymes}
\thispagestyle{plain}

\begin{table}[H]
  \centering
  \caption{Liste des acronymes utilisés dans le rapport}
  \begin{tabularx}{\textwidth}{lX}
    \toprule
    \textbf{Acronyme} & \textbf{Signification} \\
    \midrule
    AI    & Artificial Intelligence (Intelligence Artificielle) \\
    API   & Application Programming Interface \\
    GLTF  & GL Transmission Format (format de modèles 3D) \\
    HTTP  & HyperText Transfer Protocol \\
    IA    & Intelligence Artificielle \\
    JWT   & JSON Web Token \\
    MVC   & Modèle-Vue-Contrôleur \\
    ORM   & Object-Relational Mapping \\
    PFE   & Projet de Fin d'Études \\
    PMS   & Property Management System (Système de Gestion Hôtelière) \\
    QR    & Quick Response (code-barres 2D) \\
    REST  & Representational State Transfer \\
    SPA   & Single Page Application \\
    SQL   & Structured Query Language \\
    UML   & Unified Modeling Language \\
    URL   & Uniform Resource Locator \\
    USMS  & Université Sultan Moulay Slimane \\
    % Ajoutez vos acronymes ici
    \bottomrule
  \end{tabularx}
\end{table}

\newpage

% ── TABLE DES MATIÈRES ──────────────────────────────────────
\tableofcontents
\newpage

% ── LISTE DES FIGURES ───────────────────────────────────────
\listoffigures
\addcontentsline{toc}{chapter}{\listfigurename}
\newpage

% ── LISTE DES TABLEAUX ──────────────────────────────────────
\listoftables
\addcontentsline{toc}{chapter}{\listtablename}
\newpage

% ────────────────────────────────────────────────────────────
%  CORPS DU RAPPORT (numérotation arabe)
% ────────────────────────────────────────────────────────────
\pagenumbering{arabic}
\setcounter{page}{1}

% ============================================================
%  INTRODUCTION GÉNÉRALE
% ============================================================
\chapter*{Introduction Générale}
\addcontentsline{toc}{chapter}{Introduction Générale}
\markboth{Introduction Générale}{Introduction Générale}

La transformation numérique du secteur de l'hôtellerie et du tourisme s'impose aujourd'hui
comme une nécessité incontournable. Face à une clientèle de plus en plus connectée et
exigeante, les hôtels doivent moderniser leurs outils de gestion pour rester compétitifs,
améliorer l'expérience client et optimiser leur efficacité opérationnelle.

\bigskip
Les systèmes de gestion hôtelière traditionnels reposent souvent sur des processus manuels
et des outils fragmentés (tableaux Excel, registres papier, logiciels isolés) qui génèrent
des inefficacités, des erreurs et une mauvaise expérience client. La réservation en ligne,
la gestion des disponibilités en temps réel et la communication personnalisée avec les
clients sont devenus des critères essentiels de compétitivité.

\bigskip
C'est dans ce contexte que s'inscrit notre Projet de Fin d'Études~: la conception et le
développement de \textbf{Grand Palais}, une application web de gestion hôtelière moderne
et intelligente. La problématique centrale est la suivante~: \textit{comment concevoir
une plateforme numérique unifiée qui digitalise l'intégralité du parcours client —
de la réservation en ligne jusqu'au départ de l'hôtel — tout en offrant des outils
d'administration performants aux équipes internes~?}

\bigskip
Ce mémoire s'organise autour de trois chapitres structurés~:
\begin{itemize}
  \item Le \textbf{premier chapitre} présente l'étude préalable~: contexte du projet,
        analyse de l'existant et spécification des besoins fonctionnels et non fonctionnels.
  \item Le \textbf{deuxième chapitre} développe la conception du système à travers
        une modélisation UML complète (cas d'utilisation, séquences, classes).
  \item Le \textbf{troisième chapitre} expose la réalisation technique, présentant
        les choix technologiques retenus et les principales interfaces de l'application.
\end{itemize}

% ============================================================
%  CHAPITRE 1 — ÉTUDE PRÉALABLE
% ============================================================
\chapter{Étude Préalable}

\section{Introduction}

L'étude préalable constitue la première étape fondamentale du processus de développement
logiciel. Elle permet de comprendre l'environnement dans lequel le système va s'intégrer,
d'identifier les besoins réels des utilisateurs et de définir clairement les contours
fonctionnels et techniques du projet.

% ── 1.2 ─────────────────────────────────────────────────────
\section{Contexte Général du Projet}

\subsection{Présentation du contexte hôtelier}

L'industrie hôtelière mondiale connaît une mutation profonde sous l'effet de la digitalisation.
Les voyageurs réservent désormais leurs séjours principalement en ligne, comparent les
offres et s'attendent à une expérience fluide et personnalisée. Les hôtels qui ne disposent
pas d'un système de gestion numérique intégré perdent en compétitivité et en qualité de
service. Notre projet \textbf{Grand Palais} vise à répondre à cette demande.

\subsection{Présentation du projet}

Grand Palais est une plateforme web complète de gestion hôtelière (PMS — Property
Management System) développée dans le cadre du PFE. Elle centralise la gestion des
réservations, des chambres, des paiements et de la communication client en un seul
outil moderne, accessible depuis tout navigateur web.

% ── 1.3 ─────────────────────────────────────────────────────
\section{Problématique et Analyse de l'Existant}

\subsection{Problématique}

Les hôtels de taille moyenne gèrent souvent leurs réservations de manière fragmentée~:
réservations par téléphone reportées manuellement, fichiers Excel pour les disponibilités,
factures manuscrites et aucune visibilité en temps réel sur le taux d'occupation.
Cette approche génère des erreurs (doubles réservations, pertes de données), une
mauvaise expérience client et une charge administrative élevée pour le personnel.

\subsection{Étude de l'existant}

Les solutions existantes se répartissent en deux catégories~:

\begin{description}
  \item[Systèmes traditionnels] Gestion par registres papier ou tableurs Excel.
    Aucune visibilité temps réel, risque de doubles réservations, pas de traçabilité.
  \item[Logiciels propriétaires] Des solutions comme Opera PMS ou Protel existent,
    mais elles sont coûteuses, complexes à déployer et peu adaptées aux petits
    et moyens établissements.
  \item[Plateformes OTA] Des sites comme Booking.com ou Airbnb permettent la réservation
    en ligne, mais prélèvent des commissions élevées (15--25\%) et ne fournissent
    pas d'outils de gestion interne.
\end{description}

\subsection{Critique de l'existant}

L'analyse des solutions existantes révèle plusieurs lacunes majeures~:
\begin{itemize}
  \item Absence de visualisation interactive des chambres disponibles ;
  \item Aucune intégration d'intelligence artificielle pour l'assistance client ;
  \item Coût prohibitif des solutions propriétaires pour les établissements indépendants ;
  \item Pas de gestion intégrée du parcours client de bout en bout ;
  \item Manque de tableau de bord analytique en temps réel pour les gestionnaires.
\end{itemize}

% ── 1.4 ─────────────────────────────────────────────────────
\section{Solution Proposée}

\subsection{Spécification des besoins}

\subsubsection{Besoins fonctionnels}

La plateforme \textbf{Grand Palais} doit couvrir les fonctionnalités suivantes~:
\begin{itemize}
  \item Permettre aux clients de s'inscrire, se connecter et gérer leur profil ;
  \item Offrir une visualisation 3D interactive des chambres avec disponibilité
        en temps réel (chambres libres en vert, occupées en rouge) ;
  \item Gérer le cycle complet de réservation~: sélection, options de services
        additionnels, application de codes promotionnels et paiement en ligne ;
  \item Envoyer un e-mail de confirmation avec QR Code après chaque réservation ;
  \item Fournir un assistant concierge IA via WhatsApp pour répondre aux questions
        des clients sur leur séjour ;
  \item Permettre au réceptionniste de gérer les arrivées (check-in via QR Code)
        et les départs (check-out avec génération de facture) ;
  \item Offrir à l'administrateur un tableau de bord complet~: gestion des chambres,
        des utilisateurs, des statistiques et envoi de campagnes e-mail.
\end{itemize}

\subsubsection{Besoins non fonctionnels}

\begin{description}
  \item[Sécurité] Authentification via Laravel Sanctum (tokens API), contrôle d'accès
    par rôles (Spatie Permission), chiffrement des mots de passe (bcrypt).
  \item[Performance] Architecture SPA avec Vue.js~3 pour des transitions fluides,
    API REST optimisée avec pagination et indexation de la base de données.
  \item[Ergonomie] Interface premium responsive, accessible sur tous les appareils
    (ordinateur, tablette, smartphone).
  \item[Fiabilité] Paiements sécurisés via l'API Stripe~; gestion des erreurs et
    des transactions échouées.
  \item[Extensibilité] Architecture modulaire facilitant l'ajout de nouvelles
    fonctionnalités (application mobile, multi-hôtels, etc.).
\end{description}

\subsection{Identification des acteurs}

Le système Grand Palais implique trois acteurs principaux, présentés dans le
tableau~\ref{tab:acteurs}.

\begin{table}[H]
  \centering
  \caption{Acteurs du système Grand Palais et leurs responsabilités}
  \label{tab:acteurs}
  \begin{tabularx}{\textwidth}{lX}
    \toprule
    \textbf{Acteur} & \textbf{Rôle et responsabilités} \\
    \midrule
    Client          & Visiteur authentifié qui recherche des chambres, effectue des
                      réservations en ligne, paie via Stripe, et interagit avec le
                      concierge IA sur WhatsApp. \\
    Réceptionniste  & Agent de l'hôtel qui gère les arrivées (check-in par scan QR),
                      les départs (check-out avec facturation), et le suivi des
                      réservations en cours. \\
    Administrateur  & Gestionnaire qui administre le catalogue des chambres, les
                      comptes utilisateurs, les codes promotionnels, les statistiques
                      et les campagnes de communication par e-mail. \\
    \bottomrule
  \end{tabularx}
\end{table}

\section{Conclusion}

Cette étude préalable nous a permis de cerner précisément le contexte hôtelier,
d'analyser les limites des solutions existantes et de définir un ensemble cohérent
de besoins fonctionnels et non fonctionnels. Le système Grand Palais se positionne
comme une solution innovante, complète et accessible, répondant aux attentes des
trois acteurs identifiés. Ces fondations solides guident la conception UML détaillée
que nous présentons dans le chapitre suivant.

% ============================================================
%  CHAPITRE 2 — CONCEPTION
% ============================================================
\chapter{Conception}

\section{Introduction}

Ce chapitre présente la démarche de conception du système Grand Palais. Nous y
exposons la méthodologie choisie, l'analyse fonctionnelle traduite en diagrammes
UML (cas d'utilisation, séquences), ainsi que la modélisation des données qui
constituent la base de notre architecture logicielle.

% ── 2.2 ─────────────────────────────────────────────────────
\section{Méthodologie de développement}

\subsection{Choix de la méthodologie itérative}

Face à la richesse fonctionnelle du projet et à la nécessité d'intégrer des
technologies variées (3D, IA, paiement en ligne), nous avons adopté une approche
de développement itérative et incrémentale pour les raisons suivantes :
\begin{itemize}
  \item Validation progressive de chaque module avec des tests fonctionnels ;
  \item Flexibilité pour ajuster les fonctionnalités en cours de développement ;
  \item Livraison rapide de prototypes permettant de détecter les problèmes tôt ;
  \item Meilleure gestion des risques liés aux intégrations tierces (Stripe, Twilio).
\end{itemize}

\subsection{Cycles de développement}

Le développement a été organisé en quatre itérations principales :
\begin{enumerate}
  \item \textbf{Itération 1}~: Authentification, gestion des rôles et des profils ;
  \item \textbf{Itération 2}~: Catalogue des chambres, visualisation 3D et réservation ;
  \item \textbf{Itération 3}~: Paiement Stripe, QR Code, concierge IA (WhatsApp) ;
  \item \textbf{Itération 4}~: Interfaces réceptionniste, tableau de bord admin et campagnes.
\end{enumerate}

% ── 2.3 ─────────────────────────────────────────────────────
\section{Analyse fonctionnelle}

\subsection{Diagrammes de cas d'utilisation}

Pour modéliser les interactions entre les acteurs et le système, nous avons élaboré
des diagrammes de cas d'utilisation UML couvrant chaque profil utilisateur.

\subsubsection{Diagramme de cas d'utilisation — Client}

\begin{figure}[H]
  \centering
  \includegraphics[width=0.85\textwidth]{uc_client.png}
  \caption{Diagramme de cas d'utilisation — Client}
  \label{fig:uc_client}
\end{figure}

Le client peut s'inscrire, se connecter, consulter les chambres disponibles via la
maquette 3D, effectuer une réservation, payer en ligne, gérer ses réservations
et interagir avec le concierge IA sur WhatsApp.

\subsubsection{Diagramme de cas d'utilisation — Réceptionniste}

\begin{figure}[H]
  \centering
  \includegraphics[width=0.85\textwidth]{uc_reception.png}
  \caption{Diagramme de cas d'utilisation — Réceptionniste}
  \label{fig:uc_reception}
\end{figure}

Le réceptionniste peut rechercher une réservation par nom ou QR Code, valider
le check-in, gérer le check-out et générer la facture finale du client.

\subsubsection{Diagramme de cas d'utilisation — Administrateur}

\begin{figure}[H]
  \centering
  \includegraphics[width=0.85\textwidth]{uc_admin.png}
  \caption{Diagramme de cas d'utilisation — Administrateur}
  \label{fig:uc_admin}
\end{figure}

L'administrateur gère l'ensemble du back-office~: création et modification des
chambres, gestion des utilisateurs et de leurs rôles, création de codes promotionnels,
consultation des statistiques et envoi de campagnes e-mail à la base clients.

\subsection{Diagrammes de séquence}

Les diagrammes de séquence illustrent les flux d'interaction entre les acteurs et
le système pour chaque scénario métier principal.

\subsubsection{Inscription et connexion du client}

\begin{figure}[H]
  \centering
  \includegraphics[width=0.9\textwidth]{seq_auth.png}
  \caption{Diagramme de séquence — Inscription, validation et connexion client}
  \label{fig:seq_auth}
\end{figure}

Le client remplit le formulaire d'inscription~; le système vérifie les données et
envoie un e-mail de validation. Après confirmation, le client peut se connecter
et accéder à son espace personnel.

\subsubsection{Réservation 3D et paiement en ligne}

\begin{figure}[H]
  \centering
  \includegraphics[width=0.9\textwidth]{seq_reservation.png}
  \caption{Diagramme de séquence — Recherche 3D, réservation et paiement Stripe}
  \label{fig:seq_reservation}
\end{figure}

Le client sélectionne ses dates, visualise les chambres disponibles sur la maquette
3D, ajoute des services, applique un code promotionnel, puis paie via Stripe.
Un e-mail de confirmation contenant un QR Code est envoyé automatiquement.

\subsubsection{Check-in et Check-out à la réception}

\begin{figure}[H]
  \centering
  \includegraphics[width=0.9\textwidth]{seq_checkin.png}
  \caption{Diagramme de séquence — Check-in et Check-out (réceptionniste)}
  \label{fig:seq_checkin}
\end{figure}

Le réceptionniste scanne le QR Code du client pour retrouver sa réservation et valider
son arrivée. Au départ, il génère la facture finale, enregistre le check-out et
libère la chambre pour le service de nettoyage.

% ── 2.4 ─────────────────────────────────────────────────────
\subsection{Modélisation des données}

\subsubsection{Diagrammes de classe par section}

Pour construire la base de données, nous avons conçu des diagrammes de classe qui
représentent les entités métier et leurs relations.

\begin{figure}[H]
  \centering
  \includegraphics[width=0.85\textwidth]{dc_utilisateurs.png}
  \caption{Diagramme de classe — Utilisateurs et Structure}
  \label{fig:dc_users}
\end{figure}

\begin{figure}[H]
  \centering
  \includegraphics[width=0.85\textwidth]{dc_global.png}
  \caption{Diagramme de classe global du système}
  \label{fig:dc_global}
\end{figure}

\subsubsection{Description des entités principales}

Le tableau~\ref{tab:entites} décrit les entités centrales de la base de données
du système Grand Palais.

\begin{table}[H]
  \centering
  \caption{Entités principales de la base de données Grand Palais}
  \label{tab:entites}
  \begin{tabularx}{\textwidth}{lX}
    \toprule
    \textbf{Entité (Table)} & \textbf{Description} \\
    \midrule
    \texttt{users}        & Comptes de tous les utilisateurs (clients, réceptionnistes, administrateurs)
                            avec authentification, rôles, photo de profil et
                            informations personnelles (téléphone, adresse, pays). \\
    \texttt{rooms}        & Catalogue des chambres avec numéro, étage, surface, vue, statut
                            (disponible, occupée, maintenance) et prix. \\
    \texttt{room\_types}  & Types de chambres (Standard, Deluxe, Suite) avec description,
                            capacité et tarif de base. \\
    \texttt{reservations} & Réservations avec dates de séjour, statut (pending, confirmed,
                            checked\_in, checked\_out, cancelled), prix total,
                            remise appliquée et QR Code associé. \\
    \texttt{payments}     & Transactions de paiement liées à une réservation avec montant,
                            devise, méthode (carte Stripe, sur place) et statut. \\
    \texttt{promotions}   & Codes promotionnels avec pourcentage ou montant de réduction,
                            dates de validité et nombre d'utilisations. \\
    \texttt{services}     & Services additionnels proposés à la réservation (petit-déjeuner,
                            spa, transfert aéroport, etc.) avec tarification. \\
    \bottomrule
  \end{tabularx}
\end{table}

\section{Conclusion}

Dans ce chapitre, nous avons présenté la conception complète du système à travers une
analyse fonctionnelle rigoureuse et une modélisation UML structurée. Cette conception
détaillée prépare efficacement la phase de réalisation technique.

% ============================================================
%  CHAPITRE 3 — RÉALISATION
% ============================================================
\chapter{Réalisation}

\section{Introduction}

Après avoir finalisé l'analyse fonctionnelle ainsi que la conception de cette application,
nous présentons ici les choix techniques retenus pour son développement. Ce chapitre explique
l'environnement de travail utilisé, les technologies sélectionnées et les principales
fonctionnalités à travers des captures d'écran.

% ── 3.2 ─────────────────────────────────────────────────────
\section{Architecture du système}

\subsection{Implémentation de l'architecture MVC}

Le système Grand Palais repose sur une architecture \textbf{MVC découplée} :
le backend Laravel expose une API REST et le frontend Vue.js consomme
cette API de façon indépendante.

\begin{description}
  \item[Modèle] Couche de données g\'{e}r\'{e}e par l'ORM Eloquent de Laravel avec
    MySQL comme SGBD. Chaque table de la base de données correspond
    à un modèle Eloquent avec ses relations (\texttt{hasMany}, \texttt{belongsTo}).
  \item[Vue] SPA (Single Page Application) développée avec Vue.js~3.
    Les composants Vue communiquent avec le backend via des requêtes
    HTTP (Axios) vers les endpoints de l'API REST.
  \item[Contrôleur] Contrôleurs Laravel qui reçoivent les requêtes HTTP,
    appliquent la logique métier et retournent des réponses JSON
    consommées par le frontend.
\end{description}

\begin{figure}[H]
  \centering
  \includegraphics[width=0.85\textwidth]{architecture.png}
  \caption{Architecture MVC découplée du système Grand Palais}
  \label{fig:architecture}
\end{figure}

% ── 3.3 ─────────────────────────────────────────────────────
\section{Langages et Outils de Développement}

Le développement a impliqué plusieurs langages, frameworks et outils, sélectionnés
selon leur fiabilité, leurs performances et leur adéquation au projet.

\subsection{Langages de Programmation}

\subsubsection{PHP 8.2}

PHP est le langage côté serveur utilisé pour le backend. Sa version 8.2 apporte
des fonctionnalités modernes (types enum, fibers, attributs) qui améliorent la
lisibilité et la robustesse du code. Il est exécuté via le serveur Apache de XAMPP.

\subsubsection{JavaScript (ES2023) / TypeScript}

JavaScript est le langage du frontend. Utilisé avec Vue.js~3 et le bundler Vite,
il permet de créer une interface réactive et fluide sous forme de SPA.
Les configurations Vite utilisent des modules ES natifs pour des builds ultra-rapides.

\subsection{Frameworks et Bibliothèques}

\subsubsection{Laravel 12 (Backend)}

Laravel est le framework PHP choisi pour le backend. Il fournit un routage
puissant, un ORM complet (Eloquent), une gestion des files d'attente (Queues),
un système de mails (Mailable) et une intégration native avec Stripe et Twilio.
Les packages complémentaires utilisés sont :
\begin{itemize}
  \item \textbf{Spatie Permission} : gestion des rôles et permissions (admin, receptionist, client) ;
  \item \textbf{Laravel Sanctum} : authentification par tokens API pour la SPA ;
  \item \textbf{Laravel Socialite} : connexion via comptes sociaux ;
  \item \textbf{Simple QrCode} : génération des QR Codes de confirmation ;
  \item \textbf{DomPDF} : génération de factures PDF ;
  \item \textbf{OpenAI PHP Client} : intégration du concierge IA.
\end{itemize}

\subsubsection{Vue.js 3 + Vite (Frontend)}

Vue.js~3 est le framework JavaScript progressif utilisé pour le frontend.
L'architecture en composants \texttt{.vue} (SFC — Single File Components) assure
la séparation claire entre template, logique et style. Les bibliothèques
complémentaires intégrées sont :
\begin{itemize}
  \item \textbf{Vue Router 4} : navigation entre les vues de la SPA ;
  \item \textbf{Pinia} : gestion d'état globale de l'application ;
  \item \textbf{Axios} : communication avec l'API REST Laravel ;
  \item \textbf{Three.js} : rendu 3D interactif de la maquette de l'hôtel (GLTF) ;
  \item \textbf{Stripe.js} : intégration du formulaire de paiement sécurisé ;
  \item \textbf{Chart.js / vue-chartjs} : graphiques statistiques du tableau de bord ;
  \item \textbf{SweetAlert2 / vue3-toastify} : notifications et confirmations.
\end{itemize}

\subsection{Gestion des Données}

\subsubsection{MySQL}

MySQL est le SGBD relationnel utilisé pour persister l'ensemble des données
du système. Il est hébergé localement via XAMPP. La base de données contient
neuf tables principales liées par des clés étrangères avec intégrité référentielle.

\subsubsection{Eloquent ORM}

Eloquent est l'ORM natif de Laravel. Il traduit les opérations SQL en méthodes
PHP orientées objet, simplifiant les requêtes complexes grâce aux relations
(\texttt{hasMany}, \texttt{belongsTo}, \texttt{belongsToMany}) et aux scopes.

\subsection{Sécurité et Authentification}

L'authentification est gérée par \textbf{Laravel Sanctum} qui émet des tokens
d'API sécurisés stockés côté client. Le contrôle d'accès par rôles est assuré
par \textbf{Spatie Permission} : chaque route API est protégée par un middleware
vérifiant le rôle de l'utilisateur connecté (\texttt{admin}, \texttt{receptionist},
\texttt{client}). Les mots de passe sont chiffrés avec \textbf{bcrypt}.

\subsection{Outils de Développement}

\begin{description}
  \item[IDE] Visual Studio Code — environnement de développement principal
    avec extensions PHP Intelephense et Volar (Vue.js).
  \item[Modélisation UML] Enterprise Architect — pour la conception des
    diagrammes de cas d'utilisation, de séquence et de classes.
  \item[Tests API] Postman — pour tester et valider tous les endpoints REST.
  \item[Versioning] Git / GitHub — gestion du code source et collaboration.
  \item[Serveur local] XAMPP (Apache + MySQL) avec \texttt{php artisan serve}.
\end{description}

% ── 3.4 ─────────────────────────────────────────────────────
\section{Interfaces}

\subsection{Interfaces Communes}

\subsubsection{Page d'accueil}

\begin{figure}[H]
  \centering
  \includegraphics[width=0.85\textwidth]{ui_accueil.png}
  \caption{Page d'accueil de l'application Grand Palais}
  \label{fig:ui_accueil}
\end{figure}

La page d'accueil présente une interface premium avec une héro section immersive,
une présentation des types de chambres disponibles, des services de l'hôtel
et un appel à l'action vers la réservation en ligne.

\subsubsection{Page de connexion et d'inscription}

\begin{figure}[H]
  \centering
  \includegraphics[width=0.85\textwidth]{ui_connexion.png}
  \caption{Interface de connexion et d'inscription}
  \label{fig:ui_connexion}
\end{figure}

Le client saisit ses identifiants pour se connecter. Lors de l'inscription,
un e-mail de vérification est envoyé automatiquement. La connexion
via réseaux sociaux est également disponible (Laravel Socialite).

\subsubsection{Page de profil client}

\begin{figure}[H]
  \centering
  \includegraphics[width=0.85\textwidth]{ui_profil.png}
  \caption{Interface profil utilisateur}
  \label{fig:ui_profil}
\end{figure}

Le client peut modifier ses informations personnelles, changer sa photo de profil
et consulter l'historique complet de ses réservations passées.

\subsection{Interface Client --- Réservation 3D}

\subsubsection{Visualisation 3D et sélection de chambre}

\begin{figure}[H]
  \centering
  \includegraphics[width=0.85\textwidth]{ui_3d.png}
  \caption{Maquette 3D interactive --- sélection de chambre}
  \label{fig:ui_3d}
\end{figure}

Après avoir saisi ses dates de séjour, le client visualise la maquette 3D
rendue avec Three.js. Les chambres libres s'affichent en vert et les chambres
occupées en rouge. Un clic sur une chambre affiche ses détails et permet
de lancer la réservation.

\subsubsection{Paiement en ligne (Stripe)}

\begin{figure}[H]
  \centering
  \includegraphics[width=0.85\textwidth]{ui_paiement.png}
  \caption{Interface de paiement sécurisé via Stripe}
  \label{fig:ui_paiement}
\end{figure}

Le client valide sa réservation et paie via le formulaire Stripe intégré.
Un e-mail de confirmation avec QR Code est ensuite envoyé automatiquement.

\subsection{Interface Réceptionniste}

\begin{figure}[H]
  \centering
  \includegraphics[width=0.85\textwidth]{ui_dashboard_reception.png}
  \caption{Tableau de bord du réceptionniste}
  \label{fig:dashboard_reception}
\end{figure}

Le réceptionniste accède à la liste des réservations du jour, peut rechercher
un client par nom ou scanner son QR Code pour le check-in, et génère
la facture PDF lors du check-out.

\subsection{Interfaces Administrateur}

\subsubsection{Tableau de bord administratif}

\begin{figure}[H]
  \centering
  \includegraphics[width=0.85\textwidth]{ui_dashboard_admin.png}
  \caption{Tableau de bord administratif}
\end{figure}

L'administrateur consulte les statistiques en temps réel (chiffre d'affaires,
taux d'occupation) via des graphiques Chart.js.

\subsubsection{Gestion des chambres et des utilisateurs}

\begin{figure}[H]
  \centering
  \includegraphics[width=0.85\textwidth]{ui_gestion_chambres.png}
  \caption{Interface de gestion des chambres}
\end{figure}

L'administrateur crée et modifie les chambres du catalogue (type, étage,
statut) et gère les comptes utilisateurs avec attribution de rôles.

\section{Conclusion}

La réalisation de Grand Palais a nécessité la maîtrise de technologies modernes et
variées. Le choix de Laravel~12 et Vue.js~3 a permis de construire une architecture
robuste, sécurisée et évolutive. Les intégrations uniques (Three.js, Stripe, Twilio,
OpenAI) positionnent ce projet comme une solution innovante dans la gestion hôtelière.

% ============================================================
%  CONCLUSION GÉNÉRALE
% ============================================================
\chapter*{Conclusion Générale}
\addcontentsline{toc}{chapter}{Conclusion Générale}
\markboth{Conclusion Générale}{Conclusion Générale}

\section*{Synthèse du projet}

Ce projet de fin d'études nous a permis de concevoir et de développer \textbf{Grand Palais},
une application web complète de gestion hôtelière moderne. La plateforme couvre l'intégralité
du parcours client --- de la réservation en ligne avec visualisation 3D jusqu'au départ ---
tout en offrant des outils de gestion performants au réceptionniste et à l'administrateur.

\section*{Apports du système développé}

\subsection*{Valeur ajoutée pour l'hôtel}

Le système Grand Palais apporte plusieurs avantages concrets :
\begin{itemize}
  \item Disponibilités des chambres en temps réel via la maquette 3D interactive ;
  \item Automatisation du cycle de réservation avec paiement en ligne sécurisé ;
  \item Réduction de la charge administrative grâce au check-in/check-out numérique ;
  \item Communication client avancée via le concierge IA sur WhatsApp ;
  \item Tableaux de bord analytiques pour piloter l'activité en temps réel.
\end{itemize}

\subsection*{Fonctionnalités clés implémentées}

Les fonctionnalités développées couvrent l'ensemble du cycle hôtelier :
réservation 3D, paiement Stripe, confirmation QR Code, concierge IA WhatsApp,
check-in/check-out réceptionniste, facturation PDF et campagnes e-mail clients.

\section*{Défis rencontrés et solutions apportées}

Le développement a présenté plusieurs défis techniques majeurs :
\begin{description}
  \item[Visualisation 3D] Intégration de Three.js avec des modèles GLTF et
    synchronisation des couleurs des chambres avec les disponibilités en temps réel.
  \item[Paiement Stripe] Configuration des webhooks et sécurisation du flux
    de confirmation de réservation après paiement.
  \item[Concierge IA] Intégration de Twilio (WhatsApp) et de l'API OpenAI pour
    générer des réponses personnalisées basées sur le séjour du client.
  \item[Architecture SPA] Gestion de l'authentification Sanctum entre le backend
    Laravel et le frontend Vue.js découplés sur des ports différents.
\end{description}

\section*{Compétences acquises}

Ce projet nous a permis de développer et renforcer des compétences variées :
\begin{itemize}
  \item \textbf{Compétences techniques} : Laravel, Vue.js, Three.js, Stripe, Twilio, OpenAI ;
  \item \textbf{Conception logicielle} : modélisation UML, architecture REST, gestion BDD ;
  \item \textbf{Gestion de projet} : planification itérative, versioning Git ;
  \item \textbf{Travail en équipe} : collaboration sur un dépôt GitHub commun.
\end{itemize}

\section*{Perspectives d'évolution}

Plusieurs pistes d'amélioration peuvent être envisagées :
\begin{itemize}
  \item Développement d'une application mobile (iOS/Android) pour les clients ;
  \item Support multi-hôtels pour gérer plusieurs établissements ;
  \item Module de tarification dynamique selon la saison et l'occupation ;
  \item Système de fidélité avancé avec points et récompenses ;
  \item Extension du concierge IA pour traiter les commandes de services en chambre.
\end{itemize}

% ============================================================
%  ANNEXES
% ============================================================
\appendix
\chapter{Code Source}

Le code source complet du projet est disponible sur GitHub à l'adresse suivante :

\begin{center}
  \url{https://github.com/mohsinedahnaoui/new-pfe}
\end{center}

Le repository contient la structure complète du projet avec :
\begin{itemize}
  \item \textbf{Backend Laravel} (racine du dépôt) :
    \begin{itemize}
      \item \texttt{app/Http/Controllers/} — Contrôleurs REST API (chambres, réservations, paiements, utilisateurs) ;
      \item \texttt{app/Models/} — Modèles Eloquent (User, Room, Reservation, Payment, Promotion\ldots) ;
      \item \texttt{database/migrations/} — Scripts de création des tables MySQL ;
      \item \texttt{routes/api.php} — Définition de toutes les routes REST protégées par Sanctum.
    \end{itemize}
  \item \textbf{Frontend Vue.js} (\texttt{/frontend/}) :
    \begin{itemize}
      \item \texttt{src/views/} — Pages principales (accueil, réservation 3D, dashboard admin, réception) ;
      \item \texttt{src/components/} — Composants réutilisables (maquette 3D, formulaires, tableaux) ;
      \item \texttt{src/stores/} — Stores Pinia (authentification, réservations, chambres) ;
      \item \texttt{src/router/} — Configuration du routage Vue Router avec guards d'authentification.
    \end{itemize}
\end{itemize}

\chapter{Scripts SQL}

\section{Création de la base de données}

\begin{lstlisting}[language=SQL, caption={Création de la base de données Grand Palais}]
CREATE DATABASE IF NOT EXISTS grand_palais
  CHARACTER SET utf8mb4 COLLATE utf8mb4_unicode_ci;

USE grand_palais;
\end{lstlisting}

\section{Tables principales}

\begin{lstlisting}[language=SQL, caption={Table users (tous les utilisateurs)}]
CREATE TABLE users (
  id                BIGINT AUTO_INCREMENT PRIMARY KEY,
  name              VARCHAR(255) NOT NULL,
  email             VARCHAR(255) NOT NULL UNIQUE,
  email_verified_at TIMESTAMP NULL,
  password          VARCHAR(255) NULL,
  phone             VARCHAR(255) NULL,
  avatar            VARCHAR(255) NULL,
  address           VARCHAR(255) NULL,
  city              VARCHAR(255) NULL,
  country           VARCHAR(255) NULL,
  is_active         BIT(1) DEFAULT 1,
  last_login_at     TIMESTAMP NULL,
  created_at        TIMESTAMP DEFAULT CURRENT_TIMESTAMP,
  updated_at        TIMESTAMP DEFAULT CURRENT_TIMESTAMP ON UPDATE CURRENT_TIMESTAMP,
  deleted_at        TIMESTAMP NULL
);
\end{lstlisting}

\begin{lstlisting}[language=SQL, caption={Tables rooms et reservations}]
CREATE TABLE rooms (
  id              BIGINT AUTO_INCREMENT PRIMARY KEY,
  room_type_id    BIGINT NOT NULL,
  number          VARCHAR(255) NOT NULL UNIQUE,
  floor           INT NOT NULL,
  surface         DECIMAL(8,2),
  view            VARCHAR(255) NULL,
  status          ENUM('available','occupied','maintenance','blocked')
                    DEFAULT 'available',
  price_override  DECIMAL(10,2) NULL,
  created_at      TIMESTAMP DEFAULT CURRENT_TIMESTAMP,
  FOREIGN KEY (room_type_id) REFERENCES room_types(id)
);

CREATE TABLE reservations (
  id               BIGINT AUTO_INCREMENT PRIMARY KEY,
  reference        VARCHAR(255) NOT NULL UNIQUE,
  user_id          BIGINT NOT NULL,
  room_id          BIGINT NOT NULL,
  promotion_id     BIGINT NULL,
  check_in         DATE NOT NULL,
  check_out        DATE NOT NULL,
  adults           INT DEFAULT 1,
  children         INT DEFAULT 0,
  status           ENUM('pending','confirmed','checked_in',
                        'checked_out','cancelled','no_show')
                     DEFAULT 'pending',
  room_price       DECIMAL(10,2) NOT NULL,
  services_price   DECIMAL(10,2) DEFAULT 0,
  discount_amount  DECIMAL(10,2) DEFAULT 0,
  total_price      DECIMAL(10,2) NOT NULL,
  qr_code          VARCHAR(255) NULL,
  confirmed_at     TIMESTAMP NULL,
  checked_in_at    TIMESTAMP NULL,
  checked_out_at   TIMESTAMP NULL,
  created_at       TIMESTAMP DEFAULT CURRENT_TIMESTAMP,
  FOREIGN KEY (user_id) REFERENCES users(id),
  FOREIGN KEY (room_id) REFERENCES rooms(id)
);

CREATE TABLE payments (
  id             BIGINT AUTO_INCREMENT PRIMARY KEY,
  reservation_id BIGINT NOT NULL,
  user_id        BIGINT NOT NULL,
  transaction_id VARCHAR(255) UNIQUE NULL,
  amount         DECIMAL(10,2) NOT NULL,
  currency       VARCHAR(10) DEFAULT 'EUR',
  method         ENUM('card','paypal','transfer','on_arrival','wallet'),
  status         ENUM('pending','completed','failed','refunded')
                   DEFAULT 'pending',
  paid_at        TIMESTAMP NULL,
  created_at     TIMESTAMP DEFAULT CURRENT_TIMESTAMP,
  FOREIGN KEY (reservation_id) REFERENCES reservations(id),
  FOREIGN KEY (user_id) REFERENCES users(id)
);
\end{lstlisting}

% ============================================================
%  BIBLIOGRAPHIE
% ============================================================
\printbibliography[heading=bibintoc, title={Bibliographie}]

% ── Références à ajouter dans references.bib ─────────────────
% @online{laravel,
%   author  = {{Taylor Otwell}},
%   title   = {Laravel — The PHP Framework for Web Artisans},
%   url     = {https://laravel.com/docs/12.x},
%   urldate = {2025-05-01}
% }
%
% @online{vuejs,
%   author  = {{Evan You and Vue Contributors}},
%   title   = {Vue.js 3 --- The Progressive JavaScript Framework},
%   url     = {https://vuejs.org/guide/},
%   urldate = {2025-05-01}
% }
%
% @online{threejs,
%   author  = {{Mr.doob and Three.js contributors}},
%   title   = {Three.js --- JavaScript 3D Library},
%   url     = {https://threejs.org/docs/},
%   urldate = {2025-05-01}
% }
%
% @online{stripe,
%   author  = {{Stripe, Inc.}},
%   title   = {Stripe API Reference},
%   url     = {https://stripe.com/docs/api},
%   urldate = {2025-05-01}
% }
%
% @online{spatie,
%   author  = {{Spatie}},
%   title   = {Laravel Permission --- Spatie},
%   url     = {https://spatie.be/docs/laravel-permission},
%   urldate = {2025-05-01}
% }
%
% @online{twilio,
%   author  = {{Twilio, Inc.}},
%   title   = {Twilio WhatsApp API Documentation},
%   url     = {https://www.twilio.com/docs/whatsapp},
%   urldate = {2025-05-01}
% }
%
% @online{mysql,
%   title   = {MySQL 8.0 Reference Manual},
%   url     = {https://dev.mysql.com/doc/refman/8.0/en/},
%   urldate = {2025-05-01}
% }

\end{document}
% ============================================================
%  FIN DU DOCUMENT
% ============================================================
